
\chapter{INTRODUÇÃO}
\label{capitulo-introducao}
    A convergência entre práticas recreativas no ambiente digital e a aplicação de conceitos fundamentais de Ciência da Computação constitui uma maneira de extensão para difusão científica e tecnológica. Dentre as plataformas digitais interativas, o jogo eletrônico \textit{Minecraft} se destaca não apenas por sua amplitude global, mas também por sua arquitetura cliente-servidor descentralizada, que atua como porta de entrada para que usuários desenvolvam interesses em redes, administração de servidores e protocolos de comunicação.

    No contexto acadêmico do curso de Bacharelado em Ciência da Computação do Instituto Federal do Paraná (IFPR) - Câmpus Pinhais, para a disciplina de Práticas de Extensão III, esta ação foi estruturada para atender a uma demanda de jogadores, entusiastas de tecnologia e comunidades locais, a incapacidade técnica e financeira de manter servidores persistentes, de alta performance e seguros, sem incorrer em custos recorrentes, incompatíveis com a realidade da sociedade brasileira.

    A proposta será concretizada na elaboração, gravação e disponibilização de uma série de videoaulas práticas em formato de playlist pública no YouTube, cujo objetivo central é ensinar o provisionamento e a manutenção autônoma de um ambiente virtual de computação em nuvem. A metodologia adotada descomplica o uso de uma infraestrutura corporativa ao demonstrar como utilizar a camada de gratuidade permanente (\textit{Always Free}) da Oracle Cloud Infrastructure (OCI). O escopo de instruções abrange desde a criação da instância virtual até práticas recomendadas de segurança de borda, incluindo o roteamento de portas com \textit{iptables} e listas de controle de acesso (\textit{Security Lists}), gerenciamento de zonas DNS e túneis criptografados com a Cloudflare, filtragem geográfica de pacotes e a publicação de servidores web com Nginx.

    Dessa forma, o projeto articula as diretrizes da extensão universitária ao transferir competências práticas e acadêmicas de engenharia de software e infraestrutura de redes diretamente para a sociedade civil, promovendo a capacitação autônoma de administradores e transformando o lazer em experiência educacional aplicada.